\documentclass[10pt]{article}
\usepackage[T1]{fontenc}
\usepackage[utf8]{inputenc}
\usepackage[french]{babel}
\usepackage{amsmath,amssymb,mathrsfs,mathpazo}
\usepackage{geometry}
\geometry{a4paper,margin=2cm}
\usepackage{graphicx,array,tabularx,tabularray,colortbl,multirow,multicol,booktabs,enumitem,pifont,ulem,fancybox,diagbox,cancel}
\usepackage{pstricks,pstricks-add,pst-plot,pst-tree,pst-node,pst-text,pst-func,pst-math,pst-bspline,pst-3dplot,pst-eucl,pst-solides3d}
\let\pscalculate\undefined
\usepackage{pst-bezier}
\usepackage{tikz,pgfplots,tkz-tab}
\usetikzlibrary{arrows,arrows.meta,patterns,calc,positioning,trees,decorations.markings,shapes.geometric,shapes.misc,shapes.arrows,intersections,decorations.pathmorphing,decorations.pathreplacing,backgrounds,shadows,babel}
\pgfplotsset{compat=1.18}
\usepgfplotslibrary{fillbetween}
\usepackage[np]{numprint}
\usepackage{listings,makecell,pas-tableur,twoopt}
\let\barre\undefined
\usepackage[pstricks]{tablvar}
\definecolor{darkgreen}{RGB}{0,100,0}
\providecommand{\red}{\color{red}}
\providecommand{\blue}{\color{blue}}
\providecommand{\green}{\color{green}}
\providecommand{\black}{\color{black}}
\providecommand{\white}{\color{white}}
\providecommand{\gray}{\color{gray}}
\providecommand{\cyan}{\color{cyan}}
\providecommand{\magenta}{\color{magenta}}
\providecommand{\yellow}{\color{yellow}}
\providecommand{\violet}{\color{violet}}
\providecommand{\orange}{\color{orange}}
\providecommand{\euro}{\texteuro}
\providecommand{\EUR}{\texteuro}
\providecommand{\ieme}{\textsuperscript{e}}
\providecommand{\degres}{\ensuremath{^\circ}}
\providecommand{\degree}{\ensuremath{^\circ}}
\providecommand{\vv}[1]{\overrightarrow{#1}}
\providecommand{\Coord}[1]{#1}
\providecommand{\decosix}{}
\providecommand{\starredbullet}{\ensuremath{\bullet}}
\pagestyle{empty}
\setlength{\parindent}{0pt}
\usepackage{adjustbox}\DeclareUnicodeCharacter{2236}{:}\DeclareUnicodeCharacter{2190}{\ensuremath{\leftarrow}}
\providecommand{\R}{\mathbb{R}}
\providecommand{\N}{\mathbb{N}}
\providecommand{\Z}{\mathbb{Z}}
\providecommand{\Q}{\mathbb{Q}}
\providecommand{\D}{\mathbb{D}}
\providecommand{\C}{\mathbb{C}}
\providecommand{\e}{{\mathrm{e}}}
\providecommand{\vect}[1]{\overrightarrow{#1}}
\providecommand{\vectt}[1]{\overrightarrow{\mathrm{#1}}}
\providecommand{\vvt}[1]{\overrightarrow{\mathrm{#1}}}
\providecommand{\np}[2][]{{#2~#1}}
\providecommand{\vv}[1]{{\overrightarrow{#1}}}
\renewcommand{\d}{{\mathrm{d}}}
\providecommand{\degres}{{^\circ}}
\providecommand{\degre}{{^\circ}}
\providecommand{\blue}{\color{blue}}
\providecommand{\red}{\color{red}}
\providecommand{\enskip}{\quad}
\providecommand{\llbracket}{\lbrack\!\lbrack}
\providecommand{\rrbracket}{\rbrack\!\rbrack}
\providecommand{\euro}{{\text{€}}}
\providecommand{\no}{{\text{n°}}}
\providecommand{\interval}[2]{{[#1;#2]}}
\providecommand{\intervFO}[2]{{[#1;#2[}}
\providecommand{\dotfill}{\ldots}

\providecommand{\ds}{\displaystyle}
\providecommand{\textding}[1]{#1}
\providecommand{\Oij}{\(\mathrm{O};\vec{\imath},\vec{\jmath}\)}
\providecommand{\Oijk}{\(\mathrm{O};\vec{\imath},\vec{\jmath},\vec{k}\)}
\providecommand{\Ouv}{\(\mathrm{O};\vec{u},\vec{v}\)}
\providecommand{\intervFF}[2]{\ensuremath{[#1;#2]}}
\renewcommand{\euro}{\ifmmode\text{\texteuro}\else\texteuro\fi}
\DeclareUnicodeCharacter{00B0}{\ifmmode^{\circ}\else\textdegree\fi}
\DeclareUnicodeCharacter{20AC}{\ifmmode\text{\texteuro}\else\texteuro\fi}
% Énoncés : Éducation nationale. Transcription et corrigés : APMEP, Denis Vergès et contributeurs. Découpage et classement : CoopMaths, Sébastien Lozano, Éric Elter et contributeurs. Adaptation de la présentation : Maths Explorer.
% CC BY-SA : https://coopmaths.fr/www/about/
% https://coopmaths.fr/alea/static/e3c/2020/tex/e3c_2020_00_sujet25_3.tex
\begin{document}
\noindent\textit{APMEP / CoopMaths. CC BY-SA. Présentation adaptée par Maths Explorer.}\par\medskip
\medskip

En France métropolitaine, 2018 a été l'année la plus chaude d'après les relevés météorologiques. La température moyenne y a été de 14~\textcelsius{} ; elle a dépassé de 1,4~\textcelsius{} la normale de référence calculée sur la période 1981-2010. (\emph{Source : site Météo France})

\medskip

\begin{enumerate}
\item Pour modéliser la situation, on considère l'année 2018 comme l'année zéro et on suppose que cette hausse moyenne de 1,4~\textcelsius{} par an se poursuit chaque année. Pour tout entier naturel $n$, on note alors $T_n$ la température moyenne annuelle en France pour l'année 2018+$n$.
	\begin{enumerate}
		\item Quelle est la nature de la suite $\left(T_n\right)$ ainsi définie ? On donnera son premier terme et sa raison.
		\item On considère qu'au-delà d'une température moyenne de 35~\textcelsius{} les corps ne se refroidissent pas et il devient insupportable pour les humains de continuer à habiter cette région que l'on qualifie alors d'inhabitable.

Selon le modèle considéré, en quelle année la France deviendrait-elle inhabitable pour les humains ? Justifier.
	\end{enumerate}
\item À cause du réchauffement climatique, certaines régions risquent de connaître une baisse de 10\,\% par an des précipitations moyennes annuelles mesurées en millimètres (mm). Dans une région du nord de la France, les précipitations moyennes annuelles étaient de \np[mm]{673}  en 2018. On considère l'année 2018 comme l'année zéro et on suppose que cette baisse de 10\,\% par an se poursuit chaque année. Pour tout entier naturel $n$, on note $P_n$ les précipitations annuelles moyennes en mm dans cette région pour l'année $2018+n$.
	\begin{enumerate}
		\item  Quelle est la nature de la suite $\left(P_n\right)$ ainsi définie ? On donnera son premier terme et sa raison.
		\item Pour tout entier naturel $n$, exprimer $P_n$ en fonction de $n$.
		\item On donne le programme Python suivant :

\begin{lstlisting}[style=python]
def precipitations(J):
	I=673
	n=0
	while I > J:
		I = 0.9*I
		n = n+1
	return n+2018
\end{lstlisting}

L'exécution de \og precipitations(300) \fg{} renvoie la valeur 2026. Que représente cette valeur pour le problème posé ?
	\end{enumerate}
\end{enumerate}

\vspace{0,5cm}

\end{document}
